
%%FILE blueprint/src/chapters/bounding.tex
\chapter{Bounding Rotations}

\begin{lemma} \label{lem:RaRalpha}
\lean{Bounding.Rx_norm_one,Bounding.Ry_norm_one,Bounding.Rz_norm_one,Bounding.rotR_norm_one,Bounding.rotM_norm_one,Bounding.rotR'_norm_one,Bounding.rotMθ_norm_le_one,Bounding.rotMφ_norm_le_one}
\leanok
For any $\alpha, \theta,\varphi \in \R$ and $a \in \{x,y,z\}$ one has $\| R(\alpha)\| = \| R_a(\alpha)\| = \| R'(\alpha) \| =\| M(\theta, \phi)\| = 1$ and $\| M^\theta(\theta,\varphi)\|, \| M^\phi(\theta,\varphi)\| \leq 1$.
\end{lemma}
\begin{proof}
\leanok
See \cite{polyhedron.without.rupert}, Lemma 9.
\end{proof}

\begin{lemma} \label{lem:RaRa}
  \lean{Bounding.norm_rotR_sub_rotR_lt, Bounding.norm_RxL_sub_RxL_eq, Bounding.norm_RyL_sub_RyL_eq, Bounding.norm_RzL_sub_RzL_eq}
  \leanok
Let $\epsilon>0$, $|\alpha-\alphab|\leq\varepsilon$ and $a \in \{x,y,z\}$ then
$\|R_a(\alpha)-R_a({\alphab})\|=\|R(\alpha)-R(\alphab)\| < \varepsilon$.
\end{lemma}
\begin{proof}
\leanok
See \cite{polyhedron.without.rupert}, Lemma 10.
\end{proof}



\begin{lemma} \label{lem:RxRy}
\lean{Bounding.lemma12, Bounding.lemma12_lt_of_ne}
\leanok
For any $\alpha,\beta\in \mathbb{R}$ one has
\[
    \|R_x(\alpha)R_y(\beta)-\id\| \leq  \sqrt{\alpha^2+\beta^2}
\]
with strict inequality unless $\alpha = \beta = 0$.
\end{lemma}

\begin{proof}
\leanok
The composition $R_d(\alpha)R_{d'}(\beta)$ is a rotation, i.e.\ conjugate
to some $R_z(\gamma)$ by an isometry, which gives
$\|R_d(\alpha)R_{d'}(\beta)-\id\|^2 = 2(1-\cos\gamma)
= 3 - \tr(R_d(\alpha)R_{d'}(\beta))
= 3 - (\cos\alpha + \cos\beta + \cos\alpha\cos\beta)$.
Writing the right hand side as
$2(1-\cos\alpha) + 2(1-\cos\beta) - (1-\cos\alpha)(1-\cos\beta)$
and using $2(1-\cos x) \leq x^2$ (with equality only at $x = 0$)
yields the bound and the equality condition.
This is a direct route to \cite{polyhedron.without.rupert}, Lemma 12,
that avoids the Jensen-inequality argument of Lemma 11 and the
doubling induction.
\end{proof}

\begin{lemma} \label{lem:sqrt2}
\lean{Bounding.norm_M_sub_lt, Bounding.norm_X_sub_lt}
\leanok
Let $\epsilon>0$ and $|\theta-\thetab|,|\varphi-\phib| \leq \varepsilon$ then
$\|M(\theta, \phi)-M(\thetab,\phib)\|, \|X({\theta, \varphi})-X(\thetab,\phib)\| < \sqrt{2}\varepsilon.$
\end{lemma}

\begin{proof}
\uses{lem:RxRy}
\leanok
See \cite{polyhedron.without.rupert}, Lemma 13.
\end{proof}

\begin{lemma} \label{lem:XPgt0}
  \lean{Bounding.XPgt0}
  \leanok
    Let $P \in \R^3$ with $\|P\| \leq 1$. Further, let $\epsilon>0$ and $\thetab,\phib, \theta, \phi \in \R$ such that $|\thetab-\theta|, |\phib - \phi| \leq \epsilon$. If
    \(
        \langle X(\thetab,\phib),P \rangle>\sqrt{2}\varepsilon
    \)
    then
    \(
        \langle X(\theta, \phi),P \rangle>0.
    \)
\end{lemma}

\begin{proof}
\uses{lem:sqrt2}
\leanok
See \cite{polyhedron.without.rupert}, Lemma 14.
\end{proof}

\begin{lemma} \label{lem:MPgtr}
\lean{Bounding.norm_M_apply_gt}
\leanok
    Let $P \in \R^3$ with $\|P\| \leq 1$. Further, let $\epsilon, r>0$ and $\thetab,\phib, \theta, \phi \in \R$ such that $|\thetab-\theta|, |\phib - \phi| \leq \epsilon$. If
    \(
        \| M(\thetab,\phib) P \| > r + \sqrt{2}\varepsilon
    \)
    then
    \(
        \| M(\theta,\phi) P \| > r.
    \)
\end{lemma}

\begin{proof}
\uses{lem:sqrt2}
\leanok
See \cite{polyhedron.without.rupert}, Lemma 15.
Corrigendum: the triangle inquality only implies greater than *or equal to*.
\end{proof}

\begin{lemma} \label{lem:sqrt5}
  \lean{Bounding.norm_RM_sub_RM_le}
  \leanok
    Let $\epsilon>0$ and $|\theta-\thetab|,|\varphi-\phib|,|\alpha-\alphab|\leq\varepsilon$ then
    $\|R(\alpha) M(\theta, \phi)-R(\alphab)M(\thetab,\phib)\| < \sqrt{5} \varepsilon.$
\end{lemma}

\begin{proof}
\uses{lem:RxRy}
\leanok
Write $R_z$ and $R_y$ for rotations about the third and second axes.
After dropping the norm-one projection and canceling the common rotations,
it suffices to bound
\[
  \|R_z(\alpha-\alphab)R_y(\phi)
       -R_y(\phib)R_z(\theta-\thetab)\|.
\]
Insert $R_y(\Phi)$ with the ordinary midpoint $\Phi=(\phi+\phib)/2$.
On each side, the two angle changes are bounded by $\varepsilon$ and
$\varepsilon/2$. Lemma~\ref{lem:RxRy}, applied after canceling the common
rotations, bounds each difference strictly by
$\sqrt{\varepsilon^2+(\varepsilon/2)^2}=\sqrt{5}\varepsilon/2$.
If both angle changes vanish, strictness follows from $\varepsilon>0$;
otherwise it follows from the equality condition in that lemma.
The triangle inequality gives the result. The weighted intermediate angle
used in \cite{polyhedron.without.rupert}, Lemma 16, is not needed for this
uniform bound.
\end{proof}


%%FILE blueprint/src/chapters/computational.tex
\chapter{Computational Step}

\begin{definition}\label{def:nopertQ}\leanok
  We define rational approximations of the 90 noperthedron vertices
  by $\lfloor x \cdot 10^{16} \rfloor/10^{16}$.
\end{definition}

\begin{theorem}
  \label{thm:nopertQ_approximate}
  \lean{Noperthedron.KappaApprox.exact_κApprox_python}\leanok
\uses{def:nopertQ}
\tt{nopertQ} is a $\kappa$-rational approximation of the Noperthedron.
\end{theorem}
\begin{proof}\leanok
\end{proof}

\begin{theorem}
\label{thm:exists_solution_table}
\leanok
There exists a valid solution table whose zeroth row covers
\begin{align*}
\theta_1,\theta_2&\in[0,2\pi/15] \subset [0,0.42], \\
\varphi_1&\in [0,\pi] \subset [0,3.15],\\
\varphi_2&\in [0,\pi/2] \subset [0,1.58],\\
\alpha &\in [-\pi/2,\pi/2] \subset [-1.58,1.58].
\end{align*}
\end{theorem}

\begin{proof}\leanok
\uses{def:nopertQ}
By exhibiting the table and running the validity checking algorithm.
Formalized as \texttt{Noperthedron.NativeCaseAnalysis.solutionTable}
in the build-on-demand \texttt{NativeCaseAnalysis} library (checked with
\texttt{native\_decide}) and independently in the kernel-only
\texttt{KernelCaseAnalysis} library (checked with \texttt{decide +kernel},
using only the standard axioms); both are kept out of the default build
targets so that CI runs stay fast.
\end{proof}

\begin{theorem}
\label{thm:solution_global}
If a global node in the solution tree is valid, then there is no Rupert solution for its interval.
\leanok
\lean{Noperthedron.Solution.valid_global_imp_no_rupert}
\end{theorem}

\begin{proof}\leanok
\uses{thm:global_rational,thm:nopertQ_approximate}
\end{proof}

\begin{theorem}
\label{thm:solution_local}
If a local node in the solution tree is valid, then there is no Rupert solution for its interval.
\leanok
\lean{Noperthedron.Solution.valid_local_imp_no_rupert}
\end{theorem}

\begin{proof}\leanok
\uses{thm:local_rational,lemma:nopert_verts_norm_le_one,lem:congruent,thm:nopertQ_approximate,c1_c2_c3_norms}
\end{proof}


\begin{theorem}
\leanok
\lean{Noperthedron.Solution.Row.valid_imp_not_rupert_ix,
Noperthedron.Solution.valid_split_imp_no_rupert
}
\label{thm:row_valid_imp_not_rupert_ix}
\uses{def:noperthedron}
If we have a valid solution table, and in particular its $i$th row is valid,
then there is no Rupert solution of the interval of its $i$th row.
\end{theorem}

\begin{proof}
\leanok
\uses{thm:solution_global, thm:solution_local}
By strong induction on the number of rows left in the table following the $i$th. This
is because validity constrains each row to only refer to later entries.
For a split row, every point in its interval belongs to a child interval at a
larger index, so the induction hypothesis excludes a Rupert solution there.
At a leaf, apply Theorem~\ref{thm:solution_global} or Theorem~\ref{thm:solution_local}.
\end{proof}

\begin{corollary}
\leanok
\lean{Noperthedron.Solution.Row.valid_imp_not_rupert}
\label{thm:row_valid_imp_not_rupert}
\uses{def:noperthedron}
If we have a valid solution table, then there is no Rupert solution of the interval of its zeroth row.
\end{corollary}

\begin{proof}
\leanok
\uses{thm:row_valid_imp_not_rupert_ix}
Immediate special case of \cref{thm:row_valid_imp_not_rupert_ix}.
\end{proof}


%%FILE blueprint/src/chapters/global_theorem.tex
\chapter{The Global Theorem}

\begin{lemma} \label{lem:hullscalarprod}
\leanok
\lean{GlobalTheorem.hull_scalar_prod}
Suppose $V = V_1, \ldots, V_m \subseteq \R^n$ be a finite sequence of points. Suppose $\hull V$ is its convex hull.
Let $S \in \hull V$ and $w \in \R^n$ be given. then
    \[
        \langle S ,w \rangle \leq \max_{i} \langle V_i ,w\rangle
    \]
\end{lemma}

\begin{proof}
\leanok
This is a mild generalization of \cite{polyhedron.without.rupert}, Lemma 18.

Since $S \in \hull V$, we have
\[ S = \sum_{j=1}^m \lambda_j V_j \]
for some $\lambda_1,\ldots,\lambda_m \in [0,1]$ with
\[ 1 = \sum_{j=1}^m \lambda_j  \]
Therefore
\[
\langle S ,w \rangle = \left\langle  \sum_{j=1}^m \lambda_j V_j ,w \right\rangle
=  \sum_{j=1}^m \lambda_j \left\langle   V_j ,w \right\rangle
\le   \sum_{j=1}^m \lambda_j \max_{i} \langle V_i ,w\rangle
\]
\[
=  \max_{i} \langle V_i ,w\rangle \sum_{j=1}^m \lambda_j
=  \max_{i} \langle V_i ,w\rangle
\]
as required.
\end{proof}


\begin{lemma} \label{lem:leq1}
\lean{GlobalTheorem.rotation_third_partials_bounded}
\leanok
    Let $S \in \R^3$, let $w \in \R^2$ be a unit vector, and set $f(x_1,x_2,x_3) = \langle R(x_3) M(x_1,x_2)S,w \rangle$. Then for all $x_1,x_2,x_3 \in \R$ and any $i,j,k \in \{1,2,3\}$ it holds that
    \[
        \left|\frac{\dd^2 f}{\dd x_i \dd x_j}(x_1,x_2,x_3)\right|\leq \|S\|
        \qquad\text{and}\qquad
        \left|\frac{\dd^3 f}{\dd x_i \dd x_j \dd x_k}(x_1,x_2,x_3)\right|\leq \|S\|.
    \]
    (\cite{polyhedron.without.rupert} normalizes $f$ by $1/\|S\|$ and bounds the
    partials by $1$; we keep $f$ unnormalized and carry the $\|S\|$ instead,
    which composes with the remainder bound of \cref{lem:n2} at $M = \|S\| \leq 1$.)
\end{lemma}

\begin{proof}\leanok
The second-partial bound is \cite{polyhedron.without.rupert}, Lemma 19.  For
the third partials (and, in the formalization, every higher order at once) we
argue by closure instead of enumeration: the family of functions
$\pm \langle (\pm R/R')(\alpha)\, (\partial_\theta^a \partial_\phi^b
M)(\theta,\phi)\, S, w\rangle$ with $a, b \le 2$ is closed under all three
partial derivatives --- an $\alpha$-derivative steps the head around its
four-cycle $R \to R' \to -R \to -R'$, and a $\theta$- or $\phi$-derivative
bumps one grid index, folding back with a sign at the edge since
$\partial_\theta^3 = -\partial_\theta$ and
$\partial_\phi^3 = -\partial_\phi$.  Every member is bounded by $\|S\|$
because each factor has operator norm at most one, so iterating the closure
bounds all iterated partials.
\end{proof}

\begin{lemma} \label{lem:n2}
\lean{GlobalTheorem.bounded_partials_control_difference2}
\leanok
    Let $f:\R^n\to \R$ be a $C^3$-function, let
    $\varepsilon_1,\dots,\varepsilon_n \geq 0$, and let
    $x_1,\dots,x_n,y_1,\dots,y_n \in \R$ be such that
    $|x_i-y_i|\leq \varepsilon_i$ for all $i$.
    If, for some $M \in \R$,
    \(
        \left|\partial_{x_i}\partial_{x_j}\partial_{x_k}f(v)\right| \leq M
    \)
    for all $i,j,k \in \{1,\dots,n\}$ and all $v \in \R^n$, then
    \[
    |f(x)-f(y)|\leq \sum_{i=1}^n \varepsilon_i |\partial_{x_i} f(x)|
      + \frac{1}{2} \sum_{i=1}^n \sum_{j=1}^n \varepsilon_i \varepsilon_j
          |\partial_{x_i}\partial_{x_j} f(x)|
      + \frac{M}{6}\Big(\sum_{i=1}^n \varepsilon_i\Big)^3.
    \]
    (At $M = 1$ and $\varepsilon_1 = \dots = \varepsilon_n = \varepsilon$ this
    recovers the isotropic remainder $\frac{n^3}{6}\varepsilon^3$.)
\end{lemma}

\begin{proof}
\leanok
This strengthens \cite{polyhedron.without.rupert}, Lemma 20, by one Taylor
order and per-axis radii: expand $g(t) = f((1-t)x + ty)$ to second order with
Lagrange remainder, bound $|g'(0)|$ by the first sum, $|g''(0)|/2$ by the
second (the second partials are evaluated exactly at $x$, not bounded), and
$|g'''(c)| \le M (\sum_i \varepsilon_i)^3$ using the third-partial bound.
\end{proof}

\begin{lemma}
\label{lem:rotation_derivatives}
\leanok
\lean{GlobalTheorem.rotation_partials_exist, GlobalTheorem.rotation_partials_exist_outer,
 GlobalTheorem.partials_helper0, GlobalTheorem.partials_helper1, GlobalTheorem.partials_helper2,
 GlobalTheorem.partials_helper3, GlobalTheorem.partials_helper4}
The partial derivatives of all relevant rotations, projections, and inner products
used in the Global Theorem are as expected.
Specifically:
\begin{itemize}
\item \[  f^\alpha(\theta,\phi,\alpha) = \langle R'(\alpha) M(\theta, \phi) S, w \rangle \]
\item \[  f^\theta(\theta,\phi,\alpha) = \langle R(\alpha) M^\theta(\theta, \phi) S, w \rangle \]
\item \[  f^\phi(\theta,\phi,\alpha) = \langle R(\alpha) M^\phi(\theta, \phi) S, w \rangle \]
\item \[  g^\theta(\theta,\phi) = \langle  M^\theta(\theta, \phi) P, w \rangle \]
\item \[  g^\phi(\theta,\phi) = \langle  M^\phi(\theta, \phi) P, w \rangle \]
\end{itemize}
where $f(\theta,\phi,\alpha) =  \langle R(\alpha) M(\theta,\phi) S, w\rangle $
and $g(\theta,\phi) =  \langle M(\theta,\phi) P, w\rangle $.
\end{lemma}

\begin{proof}\leanok
By basic properties of derivatives.
\end{proof}

\begin{theorem}[Global Theorem] \label{thm:global}
\lean{GlobalTheorem.global_theorem}
\leanok
    Let $\PPP$ be a pointsymmetric convex polyhedron with radius $\rho =1$ and let $S \in \PPP$. Further let $\thetab_1,\phib_1,\thetab_2,\phib_2,\alphab \in \R$, let
    $\epsilon_\alpha, \epsilon_{\theta_1}, \epsilon_{\phi_1}, \epsilon_{\theta_2}, \epsilon_{\phi_2} \geq 0$
    be per-axis radii, and let $w\in\R^2$ be a unit vector. Denote $\Mib \coloneqq M(\thetab_1, \phib_1)$, $ \Miib \coloneqq M(\thetab_2, \phib_2)$ as well as $\Mib^{\theta} \coloneqq M^\theta(\thetab_1, \phib_1)$, $\Mib^{\phi} \coloneqq M^\phi(\thetab_1, \phib_1)$ and analogously for $\Miib^{\theta}, \Miib^{\phi}$. Finally set
    \begin{align*}
        G \coloneqq{}& \langle R(\alphab) \Mib S,w \rangle - \epsilon_\alpha|\langle R'(\alphab)  \Mib S,w \rangle| - \epsilon_{\theta_1}|\langle R(\alphab) \Mib^\theta S,w \rangle| - \epsilon_{\phi_1}|\langle R(\alphab) \Mib^\phi S,w \rangle|\\
        &- \frac{1}{2}\big(\epsilon_\alpha^2|\langle R(\alphab) \Mib S,w \rangle|
          + 2\epsilon_\alpha\epsilon_{\theta_1}|\langle R'(\alphab) \Mib^\theta S,w \rangle| + 2\epsilon_\alpha\epsilon_{\phi_1}|\langle R'(\alphab) \Mib^\phi S,w \rangle|\\
        &\qquad\quad + \epsilon_{\theta_1}^2|\langle R(\alphab) \Mib^{\theta\theta} S,w \rangle| + 2\epsilon_{\theta_1}\epsilon_{\phi_1}|\langle R(\alphab) \Mib^{\theta\phi} S,w \rangle| + \epsilon_{\phi_1}^2|\langle R(\alphab) \Mib^{\phi\phi} S,w \rangle|\big)\\
        &- \frac{(\epsilon_\alpha+\epsilon_{\theta_1}+\epsilon_{\phi_1})^3}{6},\\
        H_P \coloneqq{}& \langle \Miib P,w \rangle + \epsilon_{\theta_2}|\langle \Miib^\theta P,w \rangle|+\epsilon_{\phi_2}|\langle  \Miib^\varphi P,w \rangle|\\
        &+ \frac{1}{2}\big(\epsilon_{\theta_2}^2|\langle \Miib^{\theta\theta} P,w \rangle| + 2\epsilon_{\theta_2}\epsilon_{\phi_2}|\langle \Miib^{\theta\phi} P,w \rangle| + \epsilon_{\phi_2}^2|\langle \Miib^{\phi\phi} P,w \rangle|\big) + \frac{(\epsilon_{\theta_2}+\epsilon_{\phi_2})^3}{6}, \quad \text{ for } P \in \PPP.
    \end{align*}
    If $G>\max_{P\in \PPP} H_P$ then there does not exist a solution to Rupert's condition with
    $$(\theta_1,\varphi_1,\theta_2,\varphi_2,\alpha) \in U \coloneqq [\thetab_1\pm\epsilon_{\theta_1}]\times[\phib_1\pm\epsilon_{\phi_1}]\times[\thetab_2\pm\epsilon_{\theta_2}]\times[\phib_2\pm\epsilon_{\phi_2}]\times[\alphab\pm\epsilon_\alpha] \subseteq \R^5.$$

    (This is a second-order, anisotropic strengthening of
    \cite{polyhedron.without.rupert}, Theorem 17, whose penalty is first-order
    and isotropic with quadratic remainders $9\epsilon^2/2$ and $2\epsilon^2$:
    here the exact second partials at the center pose are charged with
    per-axis weights — with multiplicities from the symmetric Hessian table —
    and only the cubic Lagrange remainders
    $(\epsilon_\alpha+\epsilon_{\theta_1}+\epsilon_{\phi_1})^3/6$ and
    $(\epsilon_{\theta_2}+\epsilon_{\phi_2})^3/6$ are bounded via
    Lemma~\ref{lem:leq1}. On the diagonal
    $\epsilon_\alpha = \epsilon_{\theta_1} = \dots = \epsilon$, these are
    $\frac{3^3}{6}\epsilon^3 = \frac{9\epsilon^3}{2}$ and
    $\frac{2^3}{6}\epsilon^3 = \frac{4\epsilon^3}{3}$.)
\end{theorem}

\begin{proof}
\leanok
\uses{lem:hullscalarprod,lem:leq1,lem:n2,lem:rotation_derivatives}
See \cite{polyhedron.without.rupert}, Section~{4.2}.
\end{proof}


%%FILE blueprint/src/chapters/local_theorem.tex
\chapter{The Local Theorem}

\begin{lemma} \label{lem:pythagoras}
  \lean{Local.pythagoras}
  \leanok
    For any $P \in \mathbb{R}^3$ one has
    $\big\|M(\theta, \phi) P\big\|^2=\|P\|^2-\langle X({\theta,\varphi}),P\rangle^2$.
\end{lemma}

\begin{proof}
\leanok
See \cite{polyhedron.without.rupert}, Lemma 21.
\end{proof}

\begin{definition} \label{def:spanp}
  \lean{Local.Spanp}
  \leanok
      Given $v_1, \dots, v_n \in \R^n$ write $\mathrm{span}^+(v_1,\dots,v_n)$ for the set (simplicial cone) in $\R^n$ defined by
    \[
        \mathrm{span}^+(v_1,\dots,v_n) = \Big\{ w \in \R^n \colon \exists \lambda_1,\dots,\lambda_n > 0 \text{ s.t. } w = \sum_{i=1}^n \lambda_i v_i \Big\},
    \]
    which is the natural restriction of $\mathrm{span}(v_1,\dots,v_n)$ to positive weights.
\end{definition}

\begin{lemma} \label{lem:langles}
  \lean{Local.langles}
  \leanok
Let $V_1,V_2,V_3,Y,Z \in \R^3$ with $\|Y \|=\| Z \|$ and $Z \in \mathrm{span}^+(V_1,V_2,V_3)$. Then there exists an $i$ with
\[
    \langle V_i, Y \rangle \leq \langle V_i, Z \rangle.
\]
\end{lemma}

\begin{proof}
\leanok
Write $Z=\sum_i c_iV_i$ with all $c_i>0$. If every claimed inequality failed,
then $\langle Z,Z\rangle<\langle Z,Y\rangle$ by taking the positive weighted sum.
But Cauchy--Schwarz gives
$\langle Z,Y\rangle\leq\|Z\|\|Y\|=\|Z\|^2$, a contradiction.
Thus the assumption that $Y$ also lies in the positive cone in
\cite{polyhedron.without.rupert}, Lemma 23, is unnecessary.
\end{proof}

\begin{lemma} \label{lem:scalarprodbars}
\lean{Local.abs_sub_inner_bars_le}
\leanok
    For $A,\overline{A},B,\overline{B}\in \R^{m\times n}$ and $P_1,P_2\in \R^n$ it holds that
    \[
        |\langle AP_1,BP_2\rangle-\langle \overline{A}P_1,\overline{B}P_2\rangle|\leq \|P_1\|\cdot \|P_2\|\cdot \Big( \|A-\overline{A}\|\cdot \|\overline{B}\| +  \|\overline{A}\|\cdot \|B-\overline{B}\|+\|A-\overline{A}\|\cdot \|B-\overline{B}\|\Big).
    \]
\end{lemma}

\begin{proof}
\leanok
See \cite{polyhedron.without.rupert}, Lemma 24.
\end{proof}

\begin{lemma} \label{lem:absscalar}
\lean{Local.abs_sub_inner_le}
\leanok
    For $A,B\in \R^{m\times n}$ and $P_1,P_2\in \R^n$ one has
    $$|\langle AP_1,AP_2\rangle-\langle BP_1,BP_2\rangle|\leq \|P_1\|\cdot \|P_2\|\cdot \|A-B\|\cdot \bigg(\|A\|+\|B\| + \|A-B\|\bigg).$$
\end{lemma}

\begin{proof}
\leanok
See \cite{polyhedron.without.rupert}, Lemma 25.
\end{proof}

\begin{lemma} \label{lem:origintriangle}
\lean{Local.origin_in_triangle}
\leanok
    Let $A,B,C\in \mathbb{R}^2$ be such that $
    \langle R({\pi/2}) A,B\rangle,
    \langle R({\pi/2}) B,C\rangle,
    \langle R({\pi/2}) C,A\rangle >0$. Then the origin lies strictly in the triangle $ABC$.
\end{lemma}

\begin{proof}
\leanok
See \cite{polyhedron.without.rupert}, Lemma 26.
\end{proof}

\begin{definition}
  \label{def:eps-spanning}
  \lean{Local.Triangle.Spanning}
  \leanok
  Let $\theta, \varphi \in \mathbb{R}$, $\varepsilon > 0$, and set $M := M(\theta, \varphi)$.
  Three points $P_1, P_2, P_3 \in \mathbb{R}^3$ with $\|P_1\|, \|P_2\|, \|P_3\| \leq 1$ are
  called $\varepsilon$-spanning for $(\theta, \varphi)$ if it holds that:
\begin{align*}
    \langle R(\pi/2) M P_1,M P_{2}\rangle > 2 \epsilon(\sqrt{2} + \varepsilon),\\
    \langle R(\pi/2) M P_2,M P_{3}\rangle > 2 \epsilon(\sqrt{2} + \varepsilon),\\
    \langle R(\pi/2) M P_3,M P_{1}\rangle > 2 \epsilon(\sqrt{2} + \varepsilon).
\end{align*}
\end{definition}


\begin{lemma} \label{lem:eps-spanning}
  \uses{def:eps-spanning, def:spanp}
  \lean{Local.vecX_spanning}
  \leanok
    Let $P_1, P_2, P_3 \in \R^3$ with $\|P_1\|,\|P_2\|,\|P_3\| \leq 1$ be $\epsilon$-spanning for $(\thetab, \phib)$ and let $\theta, \phi \in \R$ such that $|\theta - \overline{\theta}|, |\phi - \overline{\phi}| \leq \epsilon$. Assume that $\langle X(\theta, \phi), P_i \rangle > 0$ for $i=1,2,3$. Then
    \[
        X(\theta, \phi) \in \spanp(P_1, P_2, P_3).
    \]
\end{lemma}

\begin{proof}
\uses{lem:scalarprodbars, lem:origintriangle, lem:sqrt2}
\leanok
See \cite{polyhedron.without.rupert}, Lemma 28.
\end{proof}

\begin{lemma} \label{lem:inCirc}
  \lean{Local.inCirc}
  \leanok
    Let $P, Q \in \R^3$ with $\|P\|, \|Q\| \leq 1$. Let $\epsilon>0$ and $\thetab_1,\phib_1,\thetab_2,\phib_2,\alphab \in \R$ with
    \[
        \|R(\alphab) M(\thetab_1, \phib_1) P - M(\thetab_2, \phib_2) Q\| \leq 2\delta.
    \]
    Finally, let $\theta_1, \phi_1, \theta_2, \phi_2, \alpha \in \R$ with $|\thetab_1-\theta_1|, |\phib_1 - \phi_1|, |\thetab_2-\theta_2|, |\phib_2-\phi_2|, |\alphab - \alpha| \leq \epsilon$. Then
    \[
        \|R(\alpha)M(\theta_1, \phi_1) P - M(\theta_2, \phi_2) Q\| < 2(\delta + \sqrt{5} \epsilon).
    \]
\end{lemma}

\begin{proof}
\uses{lem:sqrt2, lem:sqrt5}
\leanok
See \cite{polyhedron.without.rupert}, Lemma 30, which states that both points lie in the disc of radius $\delta + \sqrt{5}\epsilon$ around their midpoint at the reference pose; here the two memberships are combined into a single bound on the distance between the two points.
\end{proof}

\begin{definition}
  \label{def:LMD}
  \lean{Local.LocallyMaximallyDistant}
  \leanok
    Let $\PP \subset \R^2$ be a convex polygon and $Q \in \PP$ one of its vertices. Define $\Sect_\delta(Q) \coloneqq \Circ_{\delta}(Q) \cap \PP^\circ$ as the intersection between $\Circ_{\delta}(Q)$ and the interior of the convex hull of $\PP$.

    Moreover, $Q \in \PP$ is called \emph{$\delta$-locally maximally distant ($\delta$-LMD)} if for all $A \in \Sect_\delta(Q)$ it holds that $\|Q\| > \|A\|$.

    (This simplifies Definition 31 of \cite{polyhedron.without.rupert}, which centers the disc at an auxiliary point $\overline{Q}$ with $\|Q - \overline{Q}\| < \delta$; the only role of $\overline{Q}$ is to bound $\|A - Q\| < 2\delta$, so we center the disc at $Q$ itself, with the radius playing the role of the paper's $2\delta$.)
\end{definition}

\begin{lemma} \label{lem:LMD}
  \uses{def:LMD}
  \lean{Local.inner_ge_implies_LMD}
  \leanok
    Let $\PP$ be a convex polygon and $Q \in \PP$ be one of its vertices. Let $\delta > 0$. Assume that for some $r > 0 $ such that $\|Q\| > r$ it holds that
    \[
        \frac{\langle Q, Q - P_j \rangle}{\|Q\|\|Q - P_j\|} \geq \delta/r,
    \]
    for all other vertices $P_j \in \PP \setminus Q$. Then $Q \in \PP$ is $2\delta$-locally maximally distant.
\end{lemma}

\begin{proof}\leanok
See \cite{polyhedron.without.rupert}, Lemma 32: the disc of radius $2\delta$ around $Q$ contains the paper's disc of radius $\delta$ around any $\overline{Q}$ with $\|Q - \overline{Q}\| < \delta$ (cf.\ Definition~\ref{def:LMD}).
\end{proof}

\begin{lemma} \label{lem:coss}
  \lean{Local.coss}
  \leanok
  Let $\epsilon>0$ and $\theta,\thetab, \phi, \phib \in \R$ with $|\theta - \overline{\theta}|, |\phi - \overline{\phi}| \leq \epsilon$. Define $M = M(\theta, \phi)$ and $\overline{M} = M(\thetab, \phib)$ and let $P, Q \in \R^3$ with $\|P\|, \|Q\| \leq 1$. Assume that
    \[
        \frac{\langle \overline{M} P,\overline{M} (P-Q)\rangle - 2 \epsilon \|P-Q\| \cdot  (\sqrt{2}+\varepsilon)}{ \big(\|\overline{M} P\|+\sqrt{2} \varepsilon \big) \cdot \big(\|\overline{M}(P-Q)\|+2 \sqrt{2} \varepsilon\big)} > 0.
    \]
  Then:
    \[
        \frac{\langle {M} P,{M} (P-Q)\rangle}{\|{M} P\| \cdot \|{M}(P-Q)\|} \geq
        \frac{\langle \overline{M} P,\overline{M} (P-Q)\rangle - 2 \epsilon \|P-Q\| \cdot  (\sqrt{2}+\varepsilon)}{ \big(\|\overline{M} P\|+\sqrt{2} \varepsilon \big) \cdot \big(\|\overline{M}(P-Q)\|+2 \sqrt{2} \varepsilon\big)}.
    \]
\end{lemma}
\begin{proof}
\uses{lem:absscalar, lem:sqrt2}
\leanok
See \cite{polyhedron.without.rupert}, Lemma 33.
\end{proof}

\begin{lemma} \label{lem:congruent}
\leanok
\lean{Local.congruent_iff_sym_matrix_eq}
    Let $P_1,P_2,P_3, Q_1,Q_2,Q_3 \in \R^3$. Define the $3 \times 3$ matrices $P \coloneqq (P_1|P_2|P_3)$ and $Q \coloneqq (Q_1|Q_2|Q_3)$ and assume that $Q$ is invertible.
    Then $P_1, P_2, P_3$ and $Q_1, Q_2, Q_3$ are congruent if and only if $P^t P = Q^t Q$.
\end{lemma}

\begin{proof}
\leanok
An isometry preserves inner products, so congruence implies equality of the Gram matrices.
Conversely, since the $Q_i$ form a basis, define a linear map $L$ by $L(Q_i)=P_i$.
The equality $P^t P = Q^t Q$ says that
$\langle L(Q_i),L(Q_j)\rangle=\langle Q_i,Q_j\rangle$ for every $i,j$.
Expanding arbitrary vectors in this basis shows that $L$ preserves all inner products,
and hence is the required linear isometry.

    If $P_1, P_2, P_3$ and $Q_1, Q_2, Q_3$ are congruent then $\langle P_i,P_j\rangle = \langle LQ_i,LQ_j\rangle = \langle Q_i,Q_j\rangle$, thus $P^t P = Q^t Q$.
    For the other direction, we claim that $L \coloneqq PQ^{-1}$ is orthonormal and satisfies that $P_i = LQ_i$ for all $i=1,2,3$.
    Indeed,  $L^t L = (PQ^{-1})^t (PQ^{-1}) = (Q^t)^{-1} P^t P Q^{-1} = \mathrm{Id}$ and it holds that $LQ=PQ^{-1}Q=P$, thus $LQ_i = P_i$.
\end{proof}


\begin{theorem}[Local Theorem] \label{thm:local}
\leanok
\lean{Local.local_theorem}
    Let $\PPP$ be a polyhedron with radius $\rho=1$ and $P_1, P_2, P_3, Q_1, Q_2, Q_3 \in \PPP$ be not necessarily distinct. Assume that $P_1, P_2, P_3$ and $Q_1, Q_2, Q_3$ are congruent.

    Let $\epsilon>0$ and $\thetab_1,\phib_1,\thetab_2,\phib_2,\alphab \in \R$, then set $\Xib \coloneqq X(\thetab_1,\phib_1), \Xiib \coloneqq X(\thetab_2,\phib_2)$ as well as $\Mib \coloneqq M(\thetab_1,\phib_1), \Miib \coloneqq M(\thetab_2,\phib_2)$.
    Assume that there exist $\sigma_P, \sigma_Q \in \{0,1\}$ such that
    \[
        (-1)^{\sigma_P} \langle \Xib,P_i\rangle>\sqrt{2}\varepsilon \quad \text{and} \quad
        (-1)^{\sigma_Q} \langle \Xiib , Q_i\rangle>\sqrt{2}\varepsilon, \tag{A$_\epsilon$}
    \]
    for all $i=1,2,3$.
    Moreover, assume that $P_1,P_2,P_3$ are $\epsilon$-spanning for $(\thetab_1,\phib_1)$ and that $Q_1,Q_2,Q_3$ are $\epsilon$-spanning for $(\thetab_2,\phib_2)$.
    Finally, assume that for all $i = 1,2,3$ and any $Q_j \in \PPP \setminus Q_i$ it holds~that
    \[
        \frac{\langle \Miib Q_i,\Miib (Q_i-Q_j)\rangle - 2 \epsilon \|Q_i-Q_j\| \cdot  (\sqrt{2}+\varepsilon)}{ \big(\|\Miib Q_i\|+\sqrt{2} \varepsilon \big) \cdot \big(\|\Miib(Q_i-Q_j)\|+2 \sqrt{2} \varepsilon\big)} > \frac{\sqrt{5} \epsilon + \delta}{r}, \tag{B$_\epsilon$}
    \]
    for some $r >0$ such that $\min_{i=1,2,3}\| \Miib Q_i \| > r + \sqrt{2} \epsilon$ and for some $\delta \in \R$ with
    \[
        \delta \geq \max_{i=1,2,3}\left\|R(\alphab) \Mib P_i - \Miib Q_i\right\|/2.
    \]
    Then there exists no solution to Rupert's problem $R(\alpha) M(\theta_1,\phi_1)\PPP \subset  M(\theta_2,\phi_2)\PPP^\circ$ with
    \[
        (\theta_1, \phi_1, \theta_2, \phi_2, \alpha) \in [\thetab_1\pm\epsilon,\phib_1\pm\epsilon,\thetab_2\pm\epsilon,\phib_2\pm\epsilon,\alphab\pm\epsilon] \coloneqq U \subseteq \R^5.
    \]
\end{theorem}

\begin{proof}\leanok
\uses{lem:langles,lem:XPgt0,lem:eps-spanning,lem:MPgtr,lem:inCirc,lem:coss,lem:LMD,lem:pythagoras}
See \cite{polyhedron.without.rupert}, Theorem 36.
\end{proof}


%%FILE blueprint/src/chapters/main.tex
\chapter{Main Theorems}

\begin{theorem}
\label{thm:no_nopert_tight_pose}
\lean{Noperthedron.no_nopert_tight_pose}
\uses{def:noperthedron}
\leanok
There does not in fact exist a noperthedron Rupert solution with
\begin{align*}
\theta_1,\theta_2&\in[0,2\pi/15] \subset [0,0.42], \\
\varphi_1&\in [0,\pi] \subset [0,3.15],\\
\varphi_2&\in [0,\pi/2] \subset [0,1.58],\\
\alpha &\in [-\pi/2,\pi/2] \subset [-1.58,1.58].
\end{align*}
\end{theorem}

\begin{proof}
\leanok
\uses{thm:exists_solution_table,thm:row_valid_imp_not_rupert}
By \ref{thm:exists_solution_table}, there is a valid solution table
containing a valid row whose pose interval is a superset of
the 5-d interval above. By \ref{thm:row_valid_imp_not_rupert}, this means
there is no Rupert solution in that interval.
\end{proof}

\begin{theorem}
\label{thm:no_nopert_pose}
\lean{Noperthedron.no_nopert_pose}
\leanok
There is no 5-parameter pose that makes the noperthedron have the Rupert property.
\end{theorem}
\begin{proof}
\uses{thm:no_nopert_tight_pose, cor:rupert_tightening}
\leanok
Theorem~\ref{thm:no_nopert_tight_pose} says there is no tight pose
that makes the noperthedron Rupert. Corollary~\ref{cor:rupert_tightening}
says that this suffices for the general case.
\end{proof}


\begin{theorem}
\label{thm:no_nopert_rot_pose}
\lean{Noperthedron.no_nopert_rot_pose}
\leanok
There is no purely rotational pose that makes the noperthedron have the Rupert property.
\end{theorem}
\begin{proof}
\uses{thm:pose_of_matrix_pose, thm:no_nopert_pose}
\leanok
Suppose there were a purely rotational pose. Then convert that to an equivalent 5-parameter pose
with Theorem~\ref{thm:pose_of_matrix_pose} and
appeal to Theorem~\ref{thm:no_nopert_pose}.
\end{proof}


\begin{theorem}
\label{thm:no_nopert_matrix_pose}
\lean{Noperthedron.no_nopert_matrix_pose}
\leanok
There is no pose that makes the noperthedron have the Rupert property.
\end{theorem}
\begin{proof}
\uses{lemma:nopert_point_symmetric, thm:no_nopert_rot_pose, thm:rupert_implies_rot_rupert}
\leanok
By Theorem~\ref{thm:rupert_implies_rot_rupert}, we need only show that the noperthedron
is pointsymmetric to see that if it is Rupert, then it must be Rupert via a purely rotational pose.
But Lemma~\ref{lemma:nopert_point_symmetric} shows exactly this. And yet we know via
Theorem~\ref{thm:no_nopert_rot_pose} that the noperthedron is not rotationally Rupert, so we have
a contradiction, hence the noperthedron has no pose that makes it Rupert.
\end{proof}


\begin{theorem}
\uses{def:noperthedron}
\label{thm:nopert_not_rupert_set}
\lean{Noperthedron.nopert_not_rupert_set}
\leanok
The noperthedron is not a Rupert set.
\end{theorem}
\begin{proof}
\uses{thm:no_nopert_matrix_pose}
\leanok
By Theorem~\ref{thm:no_nopert_matrix_pose}, there is no pose that makes the noperthedron a Rupert set.
\end{proof}

\begin{theorem}
\label{thm:nopert_not_rupert}
\lean{Noperthedron.nopert_not_rupert}
\leanok
The noperthedron is not a Rupert polyhedron.
\end{theorem}

\begin{proof}
\leanok
\uses{thm:rupert_iff_rupert_set, thm:nopert_not_rupert_set}
By Theorem~\ref{thm:rupert_iff_rupert_set} it suffices to show that the convex hull of
the noperthedron vertices is not a Rupert set. But this is exactly what Theorem~\ref{thm:nopert_not_rupert_set} shows.
\end{proof}


%%FILE blueprint/src/chapters/noperthedron.tex
\chapter{The Noperthedron}

\section{Definition of the Noperthedron}

We define three points $C_1,C_2,C_3\in \mathbb{Q}^3$.
\[
    C_1\coloneqq
        \frac{1}{259375205}
        \begin{pmatrix}
        {152024884} \\ 0 \\ {210152163}
        \end{pmatrix},
\qquad
    C_2\coloneqq \frac{1}{10^{10}}
        \begin{pmatrix}
        6632738028 \\ 6106948881 \\ 3980949609
        \end{pmatrix},
\]
\[
    C_3\coloneqq
        \frac{1}{10^{10}}
        \begin{pmatrix}
        8193990033 \\ 5298215096 \\ 1230614493
        \end{pmatrix}.
\]

\begin{lemma}
\label{c1_c2_c3_norms}
\lean{Noperthedron.c1_norm_one, Noperthedron.c2_norm_bound, Noperthedron.c3_norm_bound}
\leanok
$\| C_1 \| = 1$,
${98 \over 100} < \| C_2 \| < {99 \over 100}$, and
${98 \over 100} < \| C_3 \| < {99 \over 100}$.
\end{lemma}
\begin{proof}
\leanok
Trivial arithmetic.
\end{proof}

Rotations about the $x, y, z$ axes $R_x,R_y,R_z:$  $\mathbb{R}\to \mathbb{R}^{3\times 3}$
are defined in the usual way:
    \[
            R_x(\alpha)\coloneqq
        \begin{pmatrix}
            1 & 0 & 0\\
            0 & \cos\alpha & -\sin\alpha\\
            0 & \sin\alpha & \cos\alpha
        \end{pmatrix},
        \hspace{1cm}
        R_y(\alpha)\coloneqq
        \begin{pmatrix}
            \cos\alpha & 0 & -\sin\alpha\\
            0 & 1 & 0\\
            \sin\alpha & 0 & \cos\alpha
        \end{pmatrix},
    \]
    \[
        R_z(\alpha)\coloneqq
        \begin{pmatrix}
            \cos\alpha & -\sin\alpha &0\\
            \sin\alpha & \cos\alpha &0\\
            0 & 0 & 1
        \end{pmatrix}.
    \]

We define a 30-element set $C_{30}$
\[
    \mathcal{C}_{30} \coloneqq \left\{(-1)^\ell R_z\left(\frac{2\pi k}{15}\right) \colon k=0,\dots,14; \ell=0,1\right\}.
\]
of rotations.

We write $\mathcal{C}_{30} \cdot P = \{c P \,\text{ for } \, c \in \mathcal{C}_{30}\}$ for the orbit of $P$ under the action of $\mathcal{C}_{30}$.

\begin{definition}
\label{def:noperthedron}
\lean{Noperthedron.exactVertex}
\leanok
The Noperthedron is polyhedron given by the vertex set
\[\mathcal{C}_{30} \cdot C_1 \cup \mathcal{C}_{30} \cdot C_2 \cup \mathcal{C}_{30} \cdot C_3\]
\end{definition}

\begin{lemma}
\label{lemma:nopert_verts_norm_le_one}
\lean{Noperthedron.exactVertex_norm_le_one}
\leanok
The norm of any vertex in the Noperthedron is no more than 1.
\end{lemma}
\begin{proof}
\leanok
Evident from definitions.
\end{proof}

\begin{definition}
\label{def:pointsymmetric}
\lean{PointSym}
\leanok
A set $S \subseteq \R^3$ is {\em point-symmetric} if $x \in S$ implies $-x \in S$.
\end{definition}

\begin{lemma}
\label{lemma:nopert_point_symmetric}
\lean{Noperthedron.exactPoly_point_symmetric}
\leanok
\uses{def:pointsymmetric, def:noperthedron}
The noperthedron is point-symmetric.
\end{lemma}
\begin{proof}
\leanok
\uses{lemma:pointsymmetrization_is_pointsym}
Follows from Lemma~\ref{lemma:pointsymmetrization_is_pointsym}.
\end{proof}

\section{Refined Rupert's property for the Noperthedron}

\begin{lemma} \label{lem:symmetries}
\leanok
\lean{Noperthedron.Tightening.lemma7_1,Noperthedron.Tightening.lemma7_2,Noperthedron.Tightening.lemma7_3}
Let $\PPP = \NOP$, then for all $\theta, \varphi, \alpha \in \R$, the following three identities hold (as sets):
\begin{align*}
    M({\theta+2\pi/15,\varphi})\cdot \PPP &=M(\theta, \phi) \cdot \PPP,\\
    R(\alpha+\pi)M(\theta, \phi) \cdot \PPP &=R(\alpha)M(\theta, \phi) \cdot \PPP,\\
    \begin{pmatrix}
        1&0\\
        0&-1
    \end{pmatrix}
    M(\theta, \phi) \cdot \PPP&=
    M({\theta+\pi/15,\pi-\varphi}) \cdot \PPP.
\end{align*}
\end{lemma}
\begin{proof}
\leanok
See \cite{polyhedron.without.rupert}, Lemma 7.
\end{proof}


\begin{corollary}
\label{cor:rupert_tightening}
\lean{Noperthedron.Tightening.rupert_tightening}
\leanok
If the noperthedron is Rupert, then there exists a solution with
\begin{align*}
\theta_1,\theta_2&\in[0,2\pi/15] \subset [0,0.42], \\
\varphi_1&\in [0,\pi] \subset [0,3.15],\\
\varphi_2&\in [0,\pi/2] \subset [0,1.58],\\
\alpha &\in [-\pi/2,\pi/2] \subset [-1.58,1.58].
\end{align*}
\end{corollary}

\begin{proof}
\uses{lem:symmetries}
\leanok
See \cite{polyhedron.without.rupert}, Lemma 8.
\end{proof}


%%FILE blueprint/src/chapters/prelims.tex
\chapter{Preliminaries}

TODO: This whole chapter needs organization, it's just a grab bag of miscellaneous results for now.

\section{Rupert Sets}

\begin{theorem}[Rupert Polyhedron iff Rupert Set]
\label{thm:rupert_iff_rupert_set}
\lean{rupert_iff_rupert_set}
\leanok
The following are equivalent:
\begin{itemize}
\item The convex polyhedron with vertex set $v$ is Rupert.
\item The convex closure of $v$ is a Rupert set.
\end{itemize}
\end{theorem}
\begin{proof}
\leanok
TODO: import this from the other repo
\end{proof}


\section{Poses}

TODO

\begin{theorem}
\label{thm:pose_of_matrix_pose}
\lean{pose_of_matrix_pose, converted_pose_rupert_iff}
\leanok
Given a pose with zero offset, there exists a 5-parameter pose that is equivalent to it.
\end{theorem}

\begin{proof}
By putting the pose into a canonical form as a Z rotation followed by a Y followed by a Z.
\leanok
\end{proof}

\section{Pointsymmetry and Rupertness}

\begin{theorem}
\label{thm:rupert_implies_rot_rupert}
\lean{rupert_implies_rot_rupert}
\leanok
If a set is point symmetric and convex, then it being Rupert implies
it being purely rotationally Rupert.
\end{theorem}

\begin{proof}
\leanok
TODO: informalize proof
\end{proof}




%%FILE blueprint/src/chapters/rational.tex
\chapter{Rational Versions}

\begin{definition} \label{dfn:sin_cos_approx}
\leanok
\lean{RationalApprox.sinℚ,RationalApprox.cosℚ}
    We define the two functions $\ssin, \scos: \R \to \R$ by:
    \begin{align*}
    \ssin(x) & \coloneqq x-\frac{x^3}{3}+\frac{x^5}{5!}\mp \dots +\frac{x^{25}}{25!},\\
    \scos(x) & \coloneqq 1-\frac{x^2}{2}+\frac{x^4}{4!}\mp\dots +\frac{x^{24}}{24!}.
    \end{align*}
    Further, by replacing $\sin,\cos$ with $\ssin,\scos$ we define the functions
\[
    R_{\Q}(\alpha), R'_{\Q}(\alpha), X_{\Q}(\theta, \phi), M_{\Q}(\theta, \phi), M_{\Q}^{\theta}(\theta,\varphi),M_{\Q}^{\phi}(\theta,\varphi).
\]
\end{definition}

\begin{lemma} \label{lem:sin27cos26}
\leanok
\lean{RationalApprox.sinℚ_approx,RationalApprox.cosℚ_approx}
    \[
        |\ssin(x)-\sin(x)|\leq \frac{|x|^{27}}{27!} \quad \text{and} \quad
    |\scos(x)-\cos(x)|\leq \frac{|x|^{26}}{26!}.
    \]
\uses{dfn:sin_cos_approx}
\end{lemma}
\begin{proof}
\leanok
Appeal to Taylor series bounds, using the fact that all absolute values of
higher derivatives of sine and cosine never exceed 1.
\end{proof}

\begin{lemma} \label{lem:kappa7}
\leanok
\lean{RationalApprox.sinℚ_approx',RationalApprox.cosℚ_approx'}
    For every $x\in [-4,4]$ it holds that
    \[
    |\ssin(x)-\sin(x)| \leq \frac{\kappa}{7} \quad \text{and} \quad |\scos(x)-\cos(x)|\leq \frac{\kappa}{7}.
    \]
\end{lemma}
\begin{proof}
\uses{lem:sin27cos26}
\leanok
Straightforward numerical calculation from Lemma~\ref{lem:sin27cos26}.
\end{proof}

\begin{lemma} \label{lem:A_le_deltamn}
  \lean{RationalApprox.norm_le_delta_sqrt_dims}
  \leanok
Let $A = (a_{i,j})_{1 \leq i \leq m,\ 1 \leq j \leq n} \in \R^{m \times n}$ and $\delta >0$. Assume that $|a_{i,j}| \leq \delta$. Then it holds that $\|A\| \leq \delta \sqrt{mn}.$
\end{lemma}

\begin{proof}
\leanok
    For any $v\in \R^n$ we have
    \begin{align*}
        \|Av\|^2 &=\sum_{i=1}^m \left(\sum_{j=1}^na_{i,j}v_j\right)^2 \leq \sum_{i=1}^m\left(\sum_{j=1}^n \delta |v_j|\right)^2 = \delta^2 m\left(\sum_{j=1}^n |v_j|\right)^2 \leq \delta^2 m n \|v\|^2
    \end{align*}
    using the Cauchy-Schwarz inequality. Dividing by $\|v\|$ and taking the square root proves the claim.
\end{proof}

\begin{lemma}\label{lem:dist_le_kappa}
\leanok
\lean{RationalApprox.norm_matrix_actual_approx_le_kappa}
    Let $A(x,y)$ be an $m\times n$ matrix with $1 \leq m,n\leq 3$ such that every entry is of the form $a_1(x)\cdot a_2(y)$ where
    $a_i(z)\in \{0,1,-1,\pm\sin(z),\pm\cos(z)\}.$
    Define $A_{\mathbb{Q}}(x,y)$ by replacing $\sin$ with $\ssin$ and $\cos$ with $\scos$. Then for every $x,y\in[-4,4]$ it holds that
    $\|A(x,y)-A_{\mathbb{Q}}(x,y)\|\leq\kappa$.
\end{lemma}

\begin{proof}
\leanok
\uses{lem:kappa7,lem:A_le_deltamn}
    We've replaced the assumption $a_i(z)\in \{0,1,-1,\pm\sin(z),\pm\cos(z)\}$ in \cite{polyhedron.without.rupert}'s
    Lemma~40 with $a_i(z)\in [-1,1]$.

    By assumption, for fixed $x,y$ every entry of $A(x,y)-A_{\mathbb{Q}}(x,y)$ is of the form
    $a b - \widetilde{a}\widetilde{b}$ for some $a,b\in[-1,1]$ and $|a-\widetilde{a}|,|b-\widetilde{b}|\leq \kappa/7$ by \cref{lem:kappa7}. This implies that
    \begin{align*}
      |ab-\widetilde{a}\widetilde{b}|&\leq |a b-a\widetilde{b}|+|a \widetilde{b}-\widetilde{a}\widetilde{b}|
      =|a|\cdot |b-\widetilde{b}|+|\widetilde{b}|\cdot |a-\widetilde{a}| \leq 1\cdot \kappa/7+(1+\kappa/7) \cdot \kappa/7 <\kappa/3.
    \end{align*}
    So we can apply \cref{lem:A_le_deltamn} and obtain that $\|A(x,y)-A_{\Q}(x,y)\|<\kappa/3\cdot \sqrt{3\cdot 3}=\kappa$.
\end{proof}


\begin{corollary} \label{corr:kappa1kappa}
\leanok
\lean{
RationalApprox.R_difference_norm_bounded,
RationalApprox.R'_difference_norm_bounded,
RationalApprox.M_difference_norm_bounded,
RationalApprox.Mθ_difference_norm_bounded,
RationalApprox.Mφ_difference_norm_bounded,
RationalApprox.Mθθ_difference_norm_bounded,
RationalApprox.Mθφ_difference_norm_bounded,
RationalApprox.Mφφ_difference_norm_bounded,
RationalApprox.X_difference_norm_bounded,
RationalApprox.Rℚ_norm_bounded,
RationalApprox.Mℚ_norm_bounded,
RationalApprox.R'ℚ_norm_bounded,
RationalApprox.Mθℚ_norm_bounded,
RationalApprox.Mφℚ_norm_bounded,
RationalApprox.Mθθℚ_norm_bounded,
RationalApprox.Mθφℚ_norm_bounded,
RationalApprox.Mφφℚ_norm_bounded }
    Let $\alpha,\theta,\phi\in [-4,4]$. Then it holds that
    \begin{align*}
        \|R(\alpha)-R_{\Q}(\alpha)\|, \|R'(\alpha)-R_{\Q}'(\alpha)\|,\|X(\theta,\phi)-X_{\Q}(\theta, \phi)\|, \|M(\theta, \phi)-M_{\Q}(\theta, \phi)\|, \\
        \|M^\theta(\theta,\phi)-M_{\Q}^\theta(\theta,\phi)\|,
        \|M^\phi(\theta,\phi) - M_{\Q}^\phi(\theta,\phi)\|,
        \|M^{\theta\theta}(\theta,\phi)-M_{\Q}^{\theta\theta}(\theta,\phi)\|,
        \|M^{\theta\phi}(\theta,\phi)-M_{\Q}^{\theta\phi}(\theta,\phi)\|,
        \|M^{\phi\phi}(\theta,\phi)-M_{\Q}^{\phi\phi}(\theta,\phi)\|\leq \kappa.
    \end{align*}
    Moreover,
    \[
        \|R_{\Q}(\alpha)\|, \|R'_{\Q}(\alpha)\|, \|M_{\Q}(\theta, \phi)\|, \|M_{\Q}^{\theta}(\theta,\varphi)\|, \|M_{\Q}^{\phi}(\theta,\varphi)\|,
        \|M_{\Q}^{\theta\theta}(\theta,\varphi)\|, \|M_{\Q}^{\theta\phi}(\theta,\varphi)\|, \|M_{\Q}^{\phi\phi}(\theta,\varphi)\| \leq 1+\kappa
    \]
\end{corollary}

\begin{proof}
\leanok
\uses{lem:dist_le_kappa,lem:RaRalpha}
    The first statement is a direct application of \cref{lem:dist_le_kappa} and the second statement follows immediately after using \cref{lem:RaRalpha} and the triangle inequality.
    The derivative norm bounds follow similarly, using that the operator norms of $R'$, $M^\theta$, and $M^\phi$ are at most~$1$.
\end{proof}

\begin{lemma} \label{lem:A1AnB1Bn}
\lean{RationalApprox.norm_sub_le_prod}
\leanok
    For $1 \leq i \leq n$ let $(A_i,B_i)$ be pairs of real matrices, such that for each $i$ the dimensions of $A_i$ and $B_i$ are equal. Assume moreover that the products $A_1\cdots A_n$ and $B_1 \cdots B_n$ are well defined. Finally, assume that $\|A_i-B_i\|\leq \kappa$ and let $\delta_i\geq \max(\|A_i\|,\|B_i\|,1)$. Then it holds that
    $\|A_1\cdots A_n-B_1\cdots B_n\|\leq n\kappa\cdot \delta_1\cdots \delta_n$.
\end{lemma}

\begin{proof}
\leanok
See \cite{polyhedron.without.rupert}, Lemma 42.
\end{proof}


\begin{lemma} \label{lem:boundskappa}
\leanok
\lean{
RationalApprox.bounds_kappa_M,
RationalApprox.bounds_kappa_Mθ,
RationalApprox.bounds_kappa_Mφ,
RationalApprox.bounds_kappa_RM,
RationalApprox.bounds_kappa_R'M,
RationalApprox.bounds_kappa_RMθ,
RationalApprox.bounds_kappa_RMφ,
RationalApprox.bounds_kappa_Mθθ,
RationalApprox.bounds_kappa_Mθφ,
RationalApprox.bounds_kappa_Mφφ,
RationalApprox.bounds_kappa_R'Mθ,
RationalApprox.bounds_kappa_R'Mφ,
RationalApprox.bounds_kappa_RMθθ,
RationalApprox.bounds_kappa_RMθφ,
RationalApprox.bounds_kappa_RMφφ
}
    Let $\alpha, \theta, \phi \in [-4,4]$, $P\in \R^3$ with $\|P\| \leq 1$ and let $\widetilde{P}$ be a $\kappa$-rational approximation of $P$. Set $M = M(\theta, \phi)$ and $M_{\Q} = M_{\Q}(\theta, \phi)$, $M^\theta = M^\theta(\theta, \phi)$, $M^\theta_{\Q} = M^\theta_{\Q}(\theta, \phi)$, $M^\phi = M^\phi(\theta, \phi)$, $M^\phi_{\Q} = M^\phi_{\Q}(\theta, \phi)$ as well as $R = R(\alpha)$, $R_{\Q} = R_{\Q}(\alpha)$, $R' = R'(\alpha)$, $R'_{\Q} = R'_{\Q}(\alpha)$. Finally let $w \in \R^2$ with $\|w\| = 1$. Then:
    \begin{align}
        | \langle M P, w\rangle - \langle M_{\Q} \widetilde{P}, w\rangle | & \leq 3\kappa, \label{eq:boundskappa1} \\
        | \langle M^\theta P, w\rangle - \langle M^\theta_{\Q} \widetilde{P}, w\rangle | & \leq 3\kappa,\\
        | \langle M^\phi P, w\rangle - \langle M^\phi_{\Q} \widetilde{P}, w\rangle | & \leq 3\kappa,\\
        | \langle R M P, w\rangle - \langle R_{\Q} M_{\Q} \widetilde{P}, w\rangle | & \leq 4\kappa,\label{eq:boundskappa4} \\
        | \langle R' M P, w\rangle - \langle R'_{\Q} M_{\Q} \widetilde{P}, w\rangle | & \leq 4\kappa,\\
        | \langle R M^\theta P, w\rangle - \langle R_{\Q} M^\theta_{\Q} \widetilde{P}, w\rangle | & \leq 4\kappa,\\
        | \langle R M^\phi P, w\rangle - \langle R_{\Q} M^\phi_{\Q} \widetilde{P}, w\rangle | & \leq 4\kappa,
    \end{align}
    and the same bounds hold with $M$ replaced by any of the second-derivative
    matrices $M^{\theta\theta}, M^{\theta\phi}, M^{\phi\phi}$ (at $3\kappa$
    without a rotation factor, at $4\kappa$ with $R$ in front) and with $R'$ in
    place of $R$ in front of $M^\theta$ and $M^\phi$ (at $4\kappa$).
\end{lemma}
\begin{proof}
\leanok
\uses{lem:A1AnB1Bn,corr:kappa1kappa}
See \cite{polyhedron.without.rupert}, Lemma 44.
\end{proof}


\begin{theorem}[Rational Global Theorem] \label{thm:global_rational}
\leanok
\lean{RationalApprox.GlobalTheorem.rational_global}
    Let $\PPP$ be a pointsymmetric convex polyhedron with radius $\rho =1$ and $\widetilde{\PPP}$ a $\kappa$-rational approximation. Let $\widetilde{S} \in \widetilde{\PPP}$. Further let
    $\epsilon_\alpha, \epsilon_{\theta_1}, \epsilon_{\phi_1}, \epsilon_{\theta_2}, \epsilon_{\phi_2} \in \Q_{\geq 0}$
    be per-axis radii and $\thetab_1,\phib_1,\thetab_2,\phib_2,\alphab \in \Q \cap [-4,4]$.
    Let $w\in\Q^2$ be a unit vector. Denote $\Mib \coloneqq M_{\Q}(\thetab_1, \phib_1)$, $ \Miib \coloneqq M_{\Q}(\thetab_2, \phib_2)$ as well as $\Mib^{\theta} \coloneqq M_{\Q}^\theta(\thetab_1, \phib_1)$, $\Mib^{\phi} \coloneqq M_{\Q}^\phi(\thetab_1, \phib_1)$ and analogously for $\Miib^{\theta}, \Miib^{\phi}$. Write
    $E_1 \coloneqq \epsilon_\alpha + \epsilon_{\theta_1} + \epsilon_{\phi_1}$ and
    $E_2 \coloneqq \epsilon_{\theta_2} + \epsilon_{\phi_2}$, and set
    \begin{align*}
        G^{\Q} \coloneqq{}& \langle R_{\Q}(\alphab) \Mib \widetilde{S},w \rangle - \epsilon_\alpha|\langle R_{\Q}'(\alphab)  \Mib \widetilde{S},w \rangle| - \epsilon_{\theta_1}|\langle R_{\Q}(\alphab) \Mib^\theta \widetilde{S},w \rangle| - \epsilon_{\phi_1}|\langle R_{\Q}(\alphab) \Mib^\phi \widetilde{S},w \rangle| \\
        &- \frac{1}{2}\big(\epsilon_\alpha^2|\langle R_{\Q}(\alphab) \Mib \widetilde{S},w \rangle|
          + 2\epsilon_\alpha\epsilon_{\theta_1}|\langle R_{\Q}'(\alphab) \Mib^\theta \widetilde{S},w \rangle| + 2\epsilon_\alpha\epsilon_{\phi_1}|\langle R_{\Q}'(\alphab) \Mib^\phi \widetilde{S},w \rangle|\\
        &\qquad\quad + \epsilon_{\theta_1}^2|\langle R_{\Q}(\alphab) \Mib^{\theta\theta} \widetilde{S},w \rangle| + 2\epsilon_{\theta_1}\epsilon_{\phi_1}|\langle R_{\Q}(\alphab) \Mib^{\theta\phi} \widetilde{S},w \rangle| + \epsilon_{\phi_1}^2|\langle R_{\Q}(\alphab) \Mib^{\phi\phi} \widetilde{S},w \rangle|\big)\\
        & \hspace{9cm}- \frac{E_1^3}{6} - 4\kappa \big( 1 + E_1 + \tfrac12 E_1^2\big),\\
        H^{\Q}_P \coloneqq{}& \langle \Miib P,w \rangle + \epsilon_{\theta_2}|\langle \Miib^\theta P,w \rangle|+\epsilon_{\phi_2}|\langle  \Miib^\varphi P,w \rangle|\\
        &+ \frac{1}{2}\big(\epsilon_{\theta_2}^2|\langle \Miib^{\theta\theta} P,w \rangle| + 2\epsilon_{\theta_2}\epsilon_{\phi_2}|\langle \Miib^{\theta\phi} P,w \rangle| + \epsilon_{\phi_2}^2|\langle \Miib^{\phi\phi} P,w \rangle|\big) + \frac{E_2^3}{6} + 3\kappa\big( 1+E_2+\tfrac12 E_2^2\big).
    \end{align*}
    If $G^{\Q}>\max_{P\in \widetilde{\PPP}} H^{\Q}_P$ then there does not exist a solution to Rupert's condition to $\PPP$ with
    \[
    (\theta_1,\varphi_1,\theta_2,\varphi_2,\alpha) \in [\thetab_1\pm\epsilon_{\theta_1}]\times[\phib_1\pm\epsilon_{\phi_1}]\times[\thetab_2\pm\epsilon_{\theta_2}]\times[\phib_2\pm\epsilon_{\phi_2}]\times[\alphab\pm\epsilon_\alpha].
    \]
    (On the diagonal $\epsilon_\alpha = \dots = \epsilon$ the slack terms are
    $4\kappa(1+3\epsilon+\tfrac92\epsilon^2)$ and
    $3\kappa(1+2\epsilon+2\epsilon^2)$.)
\end{theorem}

\begin{proof}
\leanok
\uses{thm:global,lem:boundskappa}
\end{proof}


\begin{definition} \label{def:ekspanning}
\leanok
\lean{Local.Triangle.κSpanning}
    Let $\theta, \phi \in \Q \cap [-4,4]$ and $M_{\Q} \coloneqq M_{\Q}(\theta, \phi)$. Three points $\widetilde{P}_1, \widetilde{P}_2, \widetilde{P}_3 \in \Q^3$ with $\|\widetilde{P}_1\|, \|\widetilde{P}_2\|, \|\widetilde{P}_3\| \leq 1+\kappa$ are called \emph{$\epsilon$-$\kappa$-spanning for $(\theta, \phi)$} if it holds that:
\begin{align*}
    \langle R(\pi/2) M_{\Q} \widetilde{P}_1,M_{\Q} \widetilde{P}_{2}\rangle > 2 \epsilon(\sqrt{2} + \epsilon) + 6\kappa,\\
    \langle R(\pi/2) M_{\Q} \widetilde{P}_2,M_{\Q} \widetilde{P}_{3}\rangle > 2 \epsilon(\sqrt{2} + \epsilon) + 6\kappa,\\
    \langle R(\pi/2) M_{\Q} \widetilde{P}_3,M_{\Q} \widetilde{P}_{1}\rangle > 2 \epsilon(\sqrt{2} + \epsilon) + 6\kappa.
\end{align*}
\end{definition}


\begin{lemma} \label{lem:ekspanningespanning}
\leanok
\lean{RationalApprox.ek_spanning_imp_e_spanning}
\uses{def:ekspanning}
    Let $P_1, P_2, P_3 \in \R^3$ with $\|P_i\| \leq 1$ and $\widetilde{P}_1, \widetilde{P}_2, \widetilde{P}_3 \in \Q^3$ be their $\kappa$-rational approximations. Assume that $\widetilde{P}_1, \widetilde{P}_2, \widetilde{P}_3$ are $\epsilon$-$\kappa$-spanning for some $\theta, \phi \in \Q \cap [-4,4]$, then $P_1, P_2, P_3$ are $\epsilon$-spanning for $\theta, \phi$.
\end{lemma}
\begin{proof}
\leanok
\uses{lem:A1AnB1Bn,lem:eps-spanning,corr:kappa1kappa}
See \cite{polyhedron.without.rupert}, Lemma 46.
\end{proof}


\begin{lemma} \label{lem:boundskappa3}
\leanok
\lean{RationalApprox.bounds_kappa3_X}
    Let $P \in \R^3$ with $\|P\|\leq 1$ and $\widetilde{P}$ a $\kappa$-rational approximation of $P$. Moreover, let $\theta, \phi \in [-4,4]$ and set $X = X(\theta, \phi)$, $X_{\Q} = X_{\Q}(\theta, \phi)$. Then
    \begin{align}
        |\langle X, P \rangle - \langle X_{\Q}, \widetilde{P} \rangle| & \leq 3 \kappa. \label{eq:boundskappa3.1}
    \end{align}
    (The inner-product and norm bounds for $M$, $M_{\Q}$ from Lemma~49 of \cite{polyhedron.without.rupert} are formalized only in the generalized forms used by \cref{lem:boundskappa4}.)
\end{lemma}
\begin{proof}
\leanok
See \cite{polyhedron.without.rupert}, Lemma 49.
\uses{lem:A1AnB1Bn}
\end{proof}


\begin{corollary} \label{lem:boundskappa4}
\leanok
\lean{RationalApprox.bounds_kappa4}
    In the setting of \cref{lem:boundskappa3}, let $\sqrt[+]{x}$ be an upper square-root function, i.e., $\sqrt{x} \leq \sqrt[+]{x}$ for all real $x \geq 0$ with rational output on rational input. Set $\|x\|_{+} \coloneqq \sqrt[+]{\|x\|^2}$. Set
    \[
        A =  \frac{\langle M P, M(P-Q)\rangle - 2 \epsilon  \|P-Q\| \cdot  (\sqrt{2}+\varepsilon)}{ \big(\| M P\|+\sqrt{2} \varepsilon \big) \cdot \big(\|M(P-Q)\|+2 \sqrt{2} \varepsilon\big)}
    \]
    as well as
    \[
        A_{\Q} =         \frac{\langle M_{\Q} \widetilde{P}, M_{\Q} (\widetilde{P}-\widetilde{Q})\rangle - 10\kappa - 2 \epsilon ( \|\widetilde{P}-\widetilde{Q}\| + 2 \kappa ) \cdot  (\sqrt{2}+\varepsilon)}{ \big(\| M_{\Q} \widetilde{P}\|_{+}+\sqrt{2} \varepsilon + 3\kappa \big) \cdot \big(\|M_{\Q}(\widetilde{P}-\widetilde{Q})\|_{+}+2 \sqrt{2} \varepsilon + 6\kappa\big)}.
    \]
    Assume that $A \geq 0$. Then it holds that $A \geq A_{\Q}$.

    (In the Lean formalization, each applied vector $M_{\Q} \widetilde{x}$ appearing in
    $A_{\Q}$ is additionally rounded down componentwise to a multiple of $10^{-13}$
    --- \texttt{Pose.rotM₂Rℚ} --- so that the checker's dot products and norms run on
    small denominators; the resulting perturbation, at most $2 \cdot 10^{-12}$ in the
    numerator and $2 \cdot 10^{-13}$ in each norm, is absorbed into the
    $10\kappa$, $3\kappa$ and $6\kappa$ terms.)
\end{corollary}
\begin{proof}
\leanok
See \cite{polyhedron.without.rupert}, Corollary 51.
\uses{lem:boundskappa3}
\end{proof}

\begin{theorem}[Rational Local Theorem]
\label{thm:local_rational}
\leanok
\lean{RationalApprox.LocalTheorem.rational_local}
\uses{def:ekspanning}
    Let $\PPP$ be a polyhedron with radius $\rho=1$ and $\widetilde{P}_i$ be a $\kappa$-rational approximation of $P_i \in \PPP$. Set $\widetilde{\PPP} = \{\widetilde{P}_i \text{ for } P_i \in \PPP \}$. Let $P_1, P_2, P_3, Q_1, Q_2, Q_3 \in \PPP$ be not necessarily distinct and assume that $P_1, P_2, P_3$ and $Q_1, Q_2, Q_3$ are congruent.
    Let $\epsilon>0$ and $\thetab_1,\phib_1,\thetab_2,\phib_2,\alphab \in \Q \cap [-4,4]$.
    Set $\Xib \coloneqq X_{\Q}(\thetab_1,\phib_1), \Xiib \coloneqq X_{\Q}(\thetab_2,\phib_2)$ as well as $\Mib \coloneqq M_{\Q}(\thetab_1,\phib_1), \Miib \coloneqq M_{\Q}(\thetab_2,\phib_2)$.
    Assume that there exist $\sigma_P, \sigma_Q \in \{0,1\}$ such that
    \[
        (-1)^{\sigma_P} \langle \Xib,\widetilde{P}_i\rangle>\sqrt{2}\varepsilon + 3\kappa \quad \text{and} \quad
        (-1)^{\sigma_Q} \langle \Xiib , \widetilde{Q}_i\rangle>\sqrt{2}\varepsilon + 3\kappa, \tag{A$^{\Q}_\epsilon$}
    \]
    for all $i=1,2,3$.
    Moreover, assume that $\widetilde{P}_1,\widetilde{P}_2,\widetilde{P}_3$ are $\epsilon$-$\kappa$-spanning for $(\thetab_1,\phib_1)$ and that $\widetilde{Q}_1,\widetilde{Q}_2,\widetilde{Q}_3$ are $\epsilon$-$\kappa$-spanning for $(\thetab_2,\phib_2)$. Let $\sqrt[+]{x}$ and $\sqrt[-]{x}$ be upper- and lower-square-root functions (bounding $\sqrt{x}$ from above/below for all real $x \geq 0$, with rational output on rational input), then set $\|Z\|_{+} \coloneqq \sqrt[+]{\|Z\|^2}$ and $\|Z\|_{-} \coloneqq \sqrt[-]{\|Z\|^2}$ for $Z \in \R^n$.
    Finally, assume that for all $i = 1,2,3$ and any $\widetilde{Q}_j \in \widetilde{\PPP} \setminus \widetilde{Q}_i$ it holds that
    \[
        \frac{\langle \Miib \widetilde{Q}_i,\Miib (\widetilde{Q}_i-\widetilde{Q}_j)\rangle - 10\kappa - 2 \epsilon ( \|\widetilde{Q}_i-\widetilde{Q}_j\|_{+} + 2 \kappa ) \cdot  (\sqrt{2}+\varepsilon)}{ \big(\|\Miib \widetilde{Q}_i\|_{+}+\sqrt{2} \varepsilon + 3\kappa \big) \cdot \big(\|\Miib(\widetilde{Q}_i-\widetilde{Q}_j)\|_{+}+2 \sqrt{2} \varepsilon + 6\kappa\big)} > \frac{\sqrt{5} \epsilon + \delta}{r}, \tag{B$^{\Q}_\epsilon$}
    \]
    for some $r >0$ such that $\min_{i=1,2,3}\| \Miib \widetilde{Q}_i \|_{-} > r + \sqrt{2} \epsilon + 3\kappa$ and for some $\delta \in \R$ with
    \[
        \delta \geq \max_{i=1,2,3}\left\|R_{\Q}(\alphab) \Mib \widetilde{P}_i-\Miib \widetilde{Q}_i\right\|_{+}/2 + 3\kappa.
    \]
    Then there exists no solution to Rupert's problem $R(\alpha) M(\theta_1,\phi_1)\PPP \subset  M(\theta_2,\phi_2)\PPP^\circ$ with
    \[
        (\theta_1, \phi_1, \theta_2, \phi_2, \alpha) \in [\thetab_1\pm\epsilon,\phib_1\pm\epsilon,\thetab_2\pm\epsilon,\phib_2\pm\epsilon,\alphab\pm\epsilon] \subseteq \R^5.
    \]

    (In the Lean formalization, the applied vector $\Miib \widetilde{Q}_i$ in the
    $r$-condition $\| \Miib \widetilde{Q}_i \|_{-} > r + \sqrt{2} \epsilon + 3\kappa$
    is additionally rounded down componentwise to a multiple of $10^{-13}$ ---
    \texttt{Pose.rotM₂Rℚ}, as in condition B$^{\Q}_\epsilon$ --- so that the
    checker's lower norm runs on small denominators; the resulting perturbation of
    at most $2 \cdot 10^{-13}$ is absorbed into the $3\kappa$ term.)
\end{theorem}

\begin{proof}
\leanok
\uses{thm:local,lem:boundskappa3,lem:boundskappa4,lem:ekspanningespanning},
\end{proof}


%%FILE blueprint/src/content.tex
\chapter{Introduction}
We follow for the most part the structure of \cite{polyhedron.without.rupert}.

\input{chapters/noperthedron.tex}
\input{chapters/bounding.tex}

\input{chapters/prelims.tex}

\input{chapters/global_theorem.tex}
\input{chapters/local_theorem.tex}
\input{chapters/rational.tex}
\input{chapters/computational.tex}

\input{chapters/main.tex}

\bibliographystyle{alpha}
\bibliography{references}

